

\SGARefTarget{sga3:ix:section:6}
\section*{6. Monomorphisms of Groups of Multiplicative Type, and
Canonical Factorization of a Homomorphism from Such a Group}
\addcontentsline{toc}{section}{6. Monomorphisms of Groups of
Multiplicative Type, and Canonical Factorization}

\medskip
\noindent\SGARefTarget{sga3:ix:lemma:6:1}\textbf{Lemma 6.1.}
\textit{Let \(S\) be a quasi-compact prescheme, and let \(G\) be a
commutative \(S\)-group prescheme, of finite presentation and
quasi-finite over \(S\). Then there is an integer \(n>0\) such that
\[
  n\cdot\operatorname{id}_G=0,
  \qquad\text{that is,}\qquad
  G={}_nG.
\]}

\emph{Proof.}
If \(S\) is the spectrum of a field \(k\), then \(G\) is finite over
\(k\), and by VII\textsubscript{A} \SGARefLink{sga3:VIIA:8.5}{8.5} it is enough to take
\[
  n=\deg(G/k).
\]
Suppose now that \(S=\operatorname{Spec}(A)\), where \(A\) is local
Artinian. Let \(k\) be the residue field of \(A\), put
\[
  G_0=G\otimes_A k,
  \qquad
  n_0=\deg(G_0/k),
\]
and distinguish two cases.

\smallskip
\SGARefTarget{sga3:ix:lemma:6:1:proof:item:a}\noindent\textit{(a) \(k\) has characteristic zero.}
Then \(G_0\) is separable over \(k\), and hence \(G\) is unramified over
\(S\). The unit section of \(G\) is therefore an open immersion, so
\[
  {}_nG=\operatorname{Ker}(n\cdot\operatorname{id}_G)
\]
is an open subscheme of \(G\). For it to equal \(G\), it is consequently
enough that it do so set-theoretically; thus one may take \(n=n_0\).

\smallskip
\SGARefTarget{sga3:ix:lemma:6:1:proof:item:b}\noindent\textit{(b) \(k\) has characteristic \(p>0\).}
Let \(\mathfrak m\) be the maximal ideal of \(A\), and put
\[
  S_r=\operatorname{Spec}(A/\mathfrak m^{r+1})
  \qquad (r\geq0).
\]
Choose an integer \(m\) such that \(\mathfrak m^{m+1}=0\). We claim that
one may take \(n=p^mn_0\). Indeed, arguing by induction on \(m\) and
putting \(n'=p^{m-1}n_0\), we may suppose that
\(n'\cdot\operatorname{id}_G\) induces the zero endomorphism on
\[
  G_{m-1}=G\times_S S_{m-1}.
\]
\footnote{Editor's note: the original has been expanded in what follows,
with the results of Expos\'e III that are used made explicit.}
Let \(E\) be the set of endomorphisms \(u\) of \(G\) having this
property, and let \(K\) be the kernel of the morphism of group functors
\[
  G\longrightarrow \prod_{S_{m-1}/S}G
\]
(cf.\ \SGARefLink{sga3:iii:remark:0-1-2}{III 0.1.2}). Then
\[
  E=\operatorname{Hom}_{S\text{-}\mathrm{gr}}(G,K);
\]
in particular, the abelian group law on \(E\) is induced by that on
\(K\). By \SGARefLink{sga3:iii:corollary:0-9}{III 0.9}, \(K\) is the \(S\)-group functor which, to every
\(g:T\to S\), assigns the \(k\)-vector space
\SGARefTarget{sga3:ix:lemma:6:1:display:01}
\[
  \operatorname{Hom}_{\mathscr O_{T_0}}
  \bigl(
    g_0^*(\Omega^1_{G_0/S_0}),
    \mathfrak m^m\mathscr O_T
  \bigr).
\]
Thus \(pu=0\) for every \(u\in E\). It follows that
\[
  pn'\cdot\operatorname{id}_G=0,
  \qquad\text{that is,}\qquad
  p^mn_0\cdot\operatorname{id}_G=0.
\]

Suppose now that \(S\) is Noetherian; this is the case to which the
general assertion reduces by the usual reduction to the Noetherian
case.
It is enough to combine the foregoing with \SGARefLink{sga3:ix:lemma:6:2}{the following lemma}.

\medskip
\noindent\SGARefTarget{sga3:ix:lemma:6:2}\textbf{Lemma 6.2.}
\textit{Let \(X\) be a quasi-finite prescheme over a Noetherian
prescheme \(S\), and let \((X_i)_{i\in I}\) be an increasing filtered
family of subpreschemes having the following property:
\[
\left\{
\begin{aligned}
 &\text{for every \(s\in S\) and every \(n\geq0\), putting}\\
 &S_{s,n}
   =\operatorname{Spec}
      \bigl(\mathscr O_{S,s}/\mathfrak m^{n+1}\bigr),
   \text{ there is an \(i\in I\) such that}\\
 &X_i\times_S S_{s,n}=X\times_S S_{s,n}.
\end{aligned}
\right.
\]
Then there is an \(i\in I\) such that \(X_i=X\).}

\emph{Proof.}
Since \(S\) is Noetherian, there is a maximal open subset \(U\) on which
\[
  X|_U=X_i|_U
\]
for all sufficiently large \(i\). We shall prove that \(U=S\). In other
words, if \(U\neq S\), we shall find an open subset \(U'\) strictly
larger than \(U\), and an \(i\), such that
\[
  X|_{U'}=X_i|_{U'}.
\]
Localizing at a maximal point \(s\) of \(S-U\), we reduce to the case in
which \(S\) is local with closed point \(s\), and
\[
  U=S-\{s\}.
\]
Indeed, put \(S'=\operatorname{Spec}(\mathscr O_{S,s})\). If
\[
  X\times_S S'=X_i\times_S S'
\]
for some \(i\),
there is an open neighborhood \(V\) of \(s\) such that
\[
  X|_V=X_i|_V.
\]
Taking \(i\) large enough that \(X_i|_U=X|_U\), we then have
\[
  X|_W=X_i|_W,
  \qquad W=U\cup V.
\]

For all sufficiently large \(i\), the equality \(X|_U=X_i|_U\) shows
that \(X_i\) is a closed subprescheme of \(X\), defined by an ideal
\(\mathscr I^{(i)}\) whose support is contained in
\[
  X_s=X_0=X\times_S S_0,
\]
where
\[
  S_n=\operatorname{Spec}(A/\mathfrak m^{n+1}),
  \qquad S=\operatorname{Spec}(A).
\]
Since \(X_0\) is quasi-finite over \(S_0\), it is a finite closed subset
of the Noetherian prescheme \(X\). Hence \(\mathscr I^{(i)}\) is a
module of finite length. There is therefore an integer \(n\geq0\) such
that
\[
  \mathscr I^{(i)}
  \cap\mathfrak m^{n+1}\mathscr O_X=0.
\]
On the other hand, by the hypothesis of the lemma, after increasing
\(i\) if necessary we may suppose that the image of
\(\mathscr I^{(i)}\) in
\[
  \mathscr O_{X_n}
  =\mathscr O_X/\mathfrak m^{n+1}\mathscr O_X
\]
is zero. It follows that \(\mathscr I^{(i)}=0\), and therefore
\(X_i=X\). \qed

\medskip
\noindent\SGARefTarget{sga3:ix:lemma:6:3}\textbf{Lemma 6.3.}
\textit{Let \(S\) be a prescheme, let \(K\) be a group prescheme over
\(S\), locally of finite presentation, and let \(s\in S\). Suppose that
\(K_s\) is quasi-finite (respectively, unramified) over \(\kappa(s)\) at
the unit element. Then there is an open neighborhood \(U\) of \(s\) such
that \(K|_U\) is locally quasi-finite (respectively, unramified) over
\(U\).}%
\begingroup
\renewcommand{\thefootnote}{\(\ast\)}%
\footnote{Cf.\ VI\textsubscript{B} \SGARefLink{concl:sga3-VIB-2.5-ii}{2.5(ii)} for a more general
statement.}%
\addtocounter{footnote}{-1}%
\endgroup

\emph{Proof.}
Let \(V\) be the set of points \(x\in K\) having the following
property: if \(t\) is the image of \(x\) in \(S\), then \(K_t\) is
quasi-finite (respectively, unramified) over \(\kappa(t)\) at \(x\);
that is, \(x\) is isolated in \(K_t\) (respectively, it is isolated and
its local ring in \(K_t\) is a separable extension of \(\kappa(t)\)).
The set \(V\) is open because \(K\) is locally of finite presentation
over \(S\).
If \(\epsilon\) denotes the unit section of \(K\), then
\(\epsilon^{-1}(V)\) is open. By hypothesis it contains \(s\), and is
therefore an open neighborhood \(U\) of \(s\). This \(U\) has the
required property. Indeed, if \(t\in U\), then \(K_t\), being a group
locally of finite type over \(\kappa(t)\), is locally quasi-finite
(respectively, unramified) over \(\kappa(t)\) because it has that
property at \(\epsilon(t)\); cf.\ VI\textsubscript{B} \SGARefLink{prop:sga3-VIB-1.3}{1.3}. \qed

Combining \SGARefLink{sga3:ix:lemma:6:1}{6.1} and
\SGARefLink{sga3:ix:lemma:6:3}{6.3}, one obtains:

\medskip
\noindent\SGARefTarget{sga3:ix:theorem:6:4}\textbf{Theorem 6.4.}
\textit{Let \(S\) be a prescheme, let \(H\) be an \(S\)-group of
multiplicative type and of finite type, and let \(K\) be a closed group
subprescheme of \(H\), of finite presentation over \(S\). Let \(s\in S\)
be such that \(K_s\) is finite. Then there are an open neighborhood
\(U\) of \(s\) and an integer \(n>0\) such that
\[
  K|_U\subset {}_nH|_U.
\]
In particular, since \({}_nH\) is finite over \(S\), the group
\(K|_U\) is finite over \(U\).}

Nakayama's lemma now gives:

\medskip
\noindent\SGARefTarget{sga3:ix:corollary:6:5}\textbf{Corollary 6.5.}
\textit{With the preceding notation, suppose that \(K_s\) is the unit
group. Then there is an open neighborhood \(U\) of \(s\) such that
\(K|_U\) is the unit group.}

\medskip
\noindent\SGARefTarget{sga3:ix:corollary:6:6}\textbf{Corollary 6.6.}
\textit{Let \(u:H\to G\) be a homomorphism of \(S\)-group preschemes,
where \(H\) is of multiplicative type and of finite type and \(G\) is
separated over \(S\). Suppose in addition that either \(S\) is locally
Noetherian or \(G\) is locally of finite presentation over
\(S\).
If \(s\in S\) and
\[
  u_s:H_s\longrightarrow G_s
\]
is a monomorphism, then there is an open neighborhood \(U\) of \(s\)
such that \(u\) induces a monomorphism
\[
  H|_U\longrightarrow G|_U.
\]}

\emph{Proof.}
Put \(K=\operatorname{Ker}(u)\). The hypothesis on \(u_s\) says that
\(K_s\) is the unit group, and the desired conclusion says that
\(K|_U\) is the unit group for a suitable \(U\). Since \(G\) is
separated over \(S\), the group \(K\) is a closed subgroup of \(H\).
If \(S\) is not assumed locally Noetherian, but \(G\) is assumed locally
of finite presentation over \(S\), then \(K\) is locally of finite
presentation over \(S\). It is therefore of finite
presentation over \(S\), since it is separated over \(S\) (as \(H\)
is) and quasi-compact over \(S\) (being closed in \(H\), which is
quasi-compact over \(S\)). Corollary
\SGARefLink{sga3:ix:corollary:6:5}{6.5} now applies. \qed

\medskip
\noindent\SGARefTarget{sga3:ix:remark:6:6:1}\textbf{Remark 6.6.1.}%
\textit{Under the hypotheses of
\SGARefLink{sga3:ix:corollary:6:6}{6.6}, note that if, in addition,
\(G\) is affine over \(S\) (respectively, of finite presentation over
\(S\)), then
\[
  H|_U\longrightarrow G|_U
\]
is even a closed immersion, as was observed in \SGARefLink{sga3:ix:corollary:2:5}{2.5} (respectively,
\SGARefLink{sga3:ix:remark:2:6}{2.6}).}

\medskip
\noindent\SGARefTarget{sga3:ix:corollary:6:7}\textbf{Corollary 6.7.}
\textit{Let \(u:H\to G\) be a homomorphism of \(S\)-group preschemes,
where \(H\) is of multiplicative type and of finite type. Suppose that
either \(S\) is locally Noetherian or \(G\) is locally of finite
presentation over \(S\). If every homomorphism on fibers
\[
  u_s:H_s\longrightarrow G_s
\]
is a monomorphism, then \(u\) is a monomorphism.}

\emph{Proof.}
The argument is the same as in
\SGARefLink{sga3:ix:corollary:6:6}{6.6}. Here the hypothesis that \(G\)
be separated over \(S\) is unnecessary in order to ensure that
\(K=\operatorname{Ker}(u)\) is closed in \(H\), since the assumption
that all the \(u_s\) are monomorphisms implies that \(K\) is
set-theoretically reduced to the unit section. \qed

\medskip
\noindent\SGARefTarget{sga3:ix:theorem:6:8}\textbf{Theorem 6.8.}
\begin{itshape}
Let \(u:H\to G\) be a homomorphism of \(S\)-group preschemes,
where \(H\) is of multiplicative type and of finite type and \(G\) is
separated over \(S\). Suppose in addition that either \(S\) is locally
Noetherian or \(G\) is of finite presentation over \(S\). Then
\[
  K=\operatorname{Ker}(u)
\]
is a subgroup of \(H\) of multiplicative type and of finite type, and
\(u\) factors as
\SGARefTarget{sga3:ix:theorem:6:8:diagram:01}
\[
\begin{tikzcd}[column sep=large]
  H
    \arrow[r,"u'"]
  & H/K
    \arrow[r,"u''"]
  & G,
\end{tikzcd}
\]
where \(H/K\) is of multiplicative type and of finite type,
\(u'\) is the canonical homomorphism and is faithfully flat (and
affine), and \(u''\) is a monomorphism.
\end{itshape}

\emph{N.B.}
As noted in \SGARefLink{sga3:ix:remark:6:6:1}{6.6.1}, the homomorphism
\(u''\) is a closed immersion if \(G\) is affine over \(S\), or if
\(G\) is of finite presentation over \(S\).

\emph{Proof.}
It is enough to prove that \(K\) is of multiplicative type; the rest
then follows from \SGARefLink{sga3:ix:proposition:2:7}{2.7} and \SGARefLink{sga3:iv:remark:5:2:6}{IV 5.2.6}.

Suppose first that \(G\) is of finite presentation over \(S\). Since
this hypothesis is stable under base change, in proving that \(K\) is
of multiplicative type we may reduce to the case in which \(H\) is
diagonalizable:
\[
  H=D_S(M),
\]
where \(M\) is a commutative group of finite type. Let \(s\in S\). The
group \(K_s\) is a closed subgroup of
\[
  H_s=D_{\kappa(s)}(M).
\]
By \SGARefLink{sga3:ix:proposition:8:1}{8.1},
it is therefore of the form \(D_{\kappa(s)}(N)\), where \(N\) is a
quotient of \(M\). Put
\[
  K'=D_S(N).
\]
Then \(K'\) is a subgroup of \(H\) of multiplicative type. Let
\[
  v':K'\longrightarrow G
\]
be the homomorphism induced by \(u\).
By construction, \(v'_s\) is the unit homomorphism. Corollary
\SGARefLink{sga3:ix:corollary:5:2}{5.2} therefore gives an open
neighborhood \(U\) of \(s\) such that
\[
  v'|_U:K'|_U\longrightarrow G|_U
\]
is the unit homomorphism. Replacing \(S\) by \(U\), we may suppose that
\(v'\) itself is the unit homomorphism, and hence that \(u\) factors as
\SGARefTarget{sga3:ix:theorem:6:8:diagram:02}
\[
\begin{tikzcd}[column sep=large]
  H
    \arrow[r,"u'"]
  & H/K'
    \arrow[r,"u''"]
  & G.
\end{tikzcd}
\]
The homomorphism \(u''_s\) is obtained from \(u_s\) by factoring through
\[
  H_s\longrightarrow H_s/\operatorname{Ker}(u_s),
\]
and is therefore a monomorphism by \SGARefLink{sga3:iv:remark:5:2:6}{IV 5.2.6}. By
\SGARefLink{sga3:ix:corollary:6:6}{6.6}, after restricting to a suitable
open neighborhood \(U\) of \(s\), the homomorphism
\[
  u''|_U:(H/K')|_U\longrightarrow G|_U
\]
is a monomorphism. Thus, after this restriction,
\[
  \operatorname{Ker}(u)
  =\operatorname{Ker}(u')
  =K',
\]
which proves that \(\operatorname{Ker}(u)\) is of multiplicative type.

The same proof applies if, instead of assuming \(G\) of finite
presentation over \(S\), one assumes \(S\) locally Noetherian, at least
when \(H\) is diagonalizable. Without this additional hypothesis on
\(H\), one must show that there is a covering base change
\[
  S'\longrightarrow S,
\]
with \(S'\) locally Noetherian, which trivializes \(H\). This will be
proved in the next expos\'e, \SGARefLink{sga3:x:corollary:4-6}{X 4.6}. \qed

\SGARefTarget{sga3:ix:section:7}
\section*{7. Algebraicity of Formal Homomorphisms into an Affine Group}
\addcontentsline{toc}{section}{7. Algebraicity of Formal Homomorphisms
into an Affine Group}

\medskip
\noindent\SGARefTarget{sga3:ix:theorem:7:1}\textbf{Theorem 7.1.}
\textit{Let \(A\) be a Noetherian ring equipped with an ideal \(I\) for
which \(A\) is separated and complete in the \(I\)-adic topology. Put
\[
  S=\operatorname{Spec}(A),
  \qquad
  S_n=\operatorname{Spec}(A/I^{n+1}),
\]
and let \(H,G\) be \(S\)-group preschemes, with \(H\) of isotrivial
multiplicative type and \(G\) affine. Then the canonical map
\SGARefTarget{sga3:ix:theorem:7:1:equation:x}
\begin{equation}
  \tag{x}
  \theta:
  \operatorname{Hom}_{S\text{-}\mathrm{gr}}(H,G)
  \longrightarrow
  \varprojlim_n
  \operatorname{Hom}_{S_n\text{-}\mathrm{gr}}(H_n,G_n)
\end{equation}
is bijective, where \(H_n,G_n\) are the \(S_n\)-groups obtained from
\(H,G\) by the base change \(S_n\to S\).}

\emph{Proof.}
Suppose first that \(H\) is diagonalizable, hence of the form
\[
  H=\operatorname{Spec}A(M),
\]
where \(B=A(M)\) is the algebra of the commutative group \(M\) with
coefficients in \(A\). We may also write
\[
  G=\operatorname{Spec}(C),
\]
where \(C\) is an \(A\)-algebra equipped with a diagonal map satisfying
the usual axioms. Homomorphisms of \(S\)-groups \(H\to G\) correspond
bijectively to homomorphisms of \(A\)-algebras
\[
  \varphi:C\longrightarrow B
\]
compatible with the diagonal maps; that is, such that, for every
\(f\in C\),
\SGARefTarget{sga3:ix:theorem:7:1:display:01}
\[
  \Delta_H(\varphi(f))
  =(\varphi\otimes\varphi)(\Delta_G(f)).
\]
There is an analogous description of the homomorphisms of \(S_n\)-groups
\(H_n\to G_n\), given by homomorphisms of \(A_n\)-algebras
\[
  \varphi_n:C_n\longrightarrow B_n,
\]
where
\[
  A_n=A/I^{n+1},
  \qquad
  B_n=B\otimes_AA_n,
  \qquad
  C_n=C\otimes_AA_n.
\]
\emph{[The source prints \(C=C\otimes_AA_n\) in the last equality;
\(C_n\) is intended.]}
Put
\SGARefTarget{sga3:ix:theorem:7:1:display:02}
\[
  \widehat B=\varprojlim_n B_n,
  \qquad
  \widehat C=\varprojlim_n C_n.
\]
Then
\[
  \widehat B=\varprojlim_n A_n(M)
\]
identifies with a submodule of the product \(A^M\), namely the submodule
of families \((a_m)_{m\in M}\) of elements of \(A\) that tend to zero
for the \(I\)-adic topology along the filter of complements of finite
subsets of \(M\). In particular, the canonical homomorphism
\[
  B\longrightarrow\widehat B
\]
is injective.\footnote{Editor's note: the argument that follows has
been expanded.}
The projective system
\[
  \theta(\varphi)=(\varphi_n)_{n\geq0}
\]
defines a homomorphism
\[
  \widehat\varphi:\widehat C\longrightarrow\widehat B
\]
such that \SGARefLink{sga3:ix:theorem:7:1:diagram:01}{the following diagram} is commutative:
\SGARefTarget{sga3:ix:theorem:7:1:diagram:01}
\[
\begin{tikzcd}[column sep=large,row sep=large]
  C
    \arrow[r,"\varphi"]
    \arrow[d]
  & B
    \arrow[d]
  \\
  \widehat C
    \arrow[r,"\widehat\varphi"]
  & \widehat B.
\end{tikzcd}
\]

We prove that \(\theta\) is surjective. Let
\((\varphi_n)_{n\geq0}\) be a projective system; we must show that it is
obtained by reduction from some \(\varphi\) in the left-hand side of
\SGARefLink{sga3:ix:theorem:7:1:equation:x}{\((x)\)}. A priori, the system \((\varphi_n)\) defines a homomorphism
on the completed algebras,
\[
  \widehat\varphi:\widehat C\longrightarrow\widehat B.
\]
It is enough to prove that the composite
\[
  \Phi:
  C\longrightarrow\widehat C
   \xrightarrow{\widehat\varphi}\widehat B
\]
sends \(C\) into \(B\). Indeed, in that case it gives a homomorphism of
\(A\)-algebras
\SGARefTarget{sga3:ix:theorem:7:1:display:phi:01}\[
  \varphi:C\longrightarrow B
\]
whose reductions are the \(\varphi_n\). \SGARefLink{sga3:ix:theorem:7:1:display:phi:01}{This homomorphism} is immediately
seen to be compatible with the diagonal maps, since the \(\varphi_n\)
are so, and since
\[
  B\otimes_AB\longrightarrow\widehat{B\otimes_AB}
\]
is injective, as follows from the same argument after replacing \(M\)
by \(M\times M\).

Passing to the completions, the diagonal homomorphism \(\Delta_H\)
defines
\SGARefTarget{sga3:ix:theorem:7:1:display:03}
\[
  \widehat\Delta_H:
  \widehat B\longrightarrow\widehat{B\otimes_AB}.
\]
Taking the projective limit of the analogous diagrams defined by the
\(\varphi_n\), one obtains the commutative diagram
\SGARefTarget{sga3:ix:theorem:7:1:diagram:02}
\[
\begin{tikzcd}[column sep=huge,row sep=large]
  C
    \arrow[r,"\Phi"]
    \arrow[d,"\Delta_G"']
  & \widehat B
    \arrow[d,"\widehat\Delta_H"]
  \\
  C\otimes_A C
    \arrow[r,"\Psi"]
  & \widehat{B\otimes_A B},
\end{tikzcd}
\]
where \(\Psi\) is the composite
\SGARefTarget{sga3:ix:theorem:7:1:display:04}
\[
  C\otimes_A C
  \xrightarrow{\ \Phi\otimes\Phi\ }
  \widehat B\otimes_A\widehat B
  \longrightarrow
  \widehat{B\otimes_A B};
\]
the last arrow is the evident canonical homomorphism. Consequently, for
every \(f\in C\), the element \(\Phi(f)\in\widehat B\) has the property
that its image under \(\widehat\Delta_H\) is a ``decomposable'' element
of \(\widehat{B\otimes_A B}\), namely an element in the image of
\(\widehat B\otimes_A\widehat B\).

Let \((e_m)_{m\in M}\) be the canonical basis of \(A(M)\), and let
\((e_{m,m'})\) be that of
\[
  A(M\times M)=A(M)\otimes_AA(M).
\]
Since
\[
  \Delta_H(e_m)=e_m\otimes e_m=e_{m,m}
  \qquad (m\in M),
\]
it remains only to apply \SGARefLink{sga3:ix:lemma:7:2}{the following lemma}.

\medskip
\noindent\SGARefTarget{sga3:ix:lemma:7:2}\textbf{Lemma 7.2.}
\begin{itshape}
Let \(A\) be a Noetherian ring, let \(M\) be a set, and let
\((a_{m,m'})\) be a family of elements of \(A\), indexed by
\(M\times M\), satisfying:
\begin{enumerate}
  \item[(i)] \SGARefTarget{sga3:ix:lemma:7:2:item:i}
  \(a_{m,m'}=0\) if \(m\neq m'\); that is, the support of the family is
  contained in the diagonal of \(M\times M\);

  \item[(ii)] \SGARefTarget{sga3:ix:lemma:7:2:item:ii}
  one has
  \SGARefTarget{sga3:ix:lemma:7:2:display:01}
  \[
    a_{m,m'}=\sum_i b_m^ic_{m'}^i,
  \]
  where the \(b^i,c^i\) are finitely many elements of \(A^M\); in other
  words, \((a_{m,m'})\) belongs to the image of the canonical
  homomorphism
  \[
    A^M\otimes_AA^M\longrightarrow A^{M\times M}.
  \]
\end{enumerate}
Then the support of \((a_{m,m'})\) is finite.
\end{itshape}

\emph{Proof.}
By \SGARefLink{sga3:ix:lemma:7:2:item:i}{(i)}, the family \((a_{m,m'})\) is determined by the elements
\[
  a_m=a_{m,m}.
\]
For \(x=(x_n)_{n\in M}\in A(M)\), put
\[
  (u\cdot x)_m=\sum_{m'}a_{m,m'}x_{m'}.
\]
Thus \((a_{m,m'})\) is viewed as the matrix of a homomorphism
\[
  u:A(M)\longrightarrow A^M.
\]
Condition \SGARefLink{sga3:ix:lemma:7:2:item:i}{(i)} gives simply
\[
  (u\cdot x)_m=a_mx_m.
\]
On the other hand, condition \SGARefLink{sga3:ix:lemma:7:2:item:ii}{(ii)} gives
\[
  u\cdot x\in\sum_i A\cdot b^i.
\]
Hence \(u(A(M))\) lies in a finitely generated \(A\)-module. If
\((e_m)_{m\in M}\) denotes the canonical basis of
\[
  A(M)\subset A^M,
\]
then all the elements \(a_me_m\) lie in a finitely generated
\(A\)-module. Since \(A\) is Noetherian, the module generated by these
elements is itself finitely generated. Because the \(e_m\) are linearly
independent, all but finitely many of the \(a_m\) are zero. This proves
the lemma, and hence Theorem \SGARefLink{sga3:ix:theorem:7:1}{7.1} when
\(H\) is diagonalizable. \qed

We now prove the general case of
\SGARefLink{sga3:ix:theorem:7:1}{7.1}, in which \(H\) is only assumed
isotrivial. There is then a finite \'etale surjective morphism
\[
  S'\longrightarrow S
\]
such that
\[
  H'=H\times_SS'
\]
is diagonalizable. We shall use only the fact that \(S'\to S\) is finite
and a covering for the faithfully flat quasi-compact topology, or simply
for the canonical topology on \((\mathbf{Sch})\). Thus ``\'etale'' could
be replaced by ``flat.''

Put
\[
  S''=S'\times_SS',
\]
and similarly introduce
\[
  S'_n
  \quad\text{and}\quad
  S''_n
  =S''\times_SS_n
  =S'_n\times_{S_n}S'_n,
\]
together with \(H',G',H'',G'',H'_n,G'_n,H''_n,G''_n\), obtained from
\(H,G\) by the evident base changes. The groups \(H'\) and \(H''\) are
diagonalizable. Moreover, \(S'\), and hence \(S''\), is affine. If
\[
  S'=\operatorname{Spec}(A'),
  \qquad
  S''=\operatorname{Spec}(A''),
\]
then \(A'\) and \(A''\) are separated and complete for the topologies
defined by \(IA'\) and \(IA''\), respectively, because both are finite
over \(A\).

Since \(S'\to S\) and \(S'_n\to S_n\) are covering morphisms, we obtain
a commutative diagram of maps of sets whose two rows are exact:
\SGARefTarget{sga3:ix:theorem:7:1:diagram:03}
\[
\begin{tikzcd}[
  column sep=small,
  row sep=large,
  cells={nodes={font=\small}}
]
  \operatorname{Hom}_{S\text{-}\mathrm{gr}}(H,G)
    \arrow[r]
    \arrow[d,"u"]
  & \operatorname{Hom}_{S'\text{-}\mathrm{gr}}(H',G')
    \arrow[r,shift left=.6ex]
    \arrow[r,shift right=.6ex]
    \arrow[d,"u'"]
  & \operatorname{Hom}_{S''\text{-}\mathrm{gr}}(H'',G'')
    \arrow[d,"u''"]
  \\
  \varprojlim_n
    \operatorname{Hom}_{S_n\text{-}\mathrm{gr}}(H_n,G_n)
    \arrow[r]
  & \varprojlim_n
    \operatorname{Hom}_{S'_n\text{-}\mathrm{gr}}(H'_n,G'_n)
    \arrow[r,shift left=.6ex]
    \arrow[r,shift right=.6ex]
  & \varprojlim_n
    \operatorname{Hom}_{S''_n\text{-}\mathrm{gr}}(H''_n,G''_n).
\end{tikzcd}
\]
The second row is exact because it is the projective limit of the exact
diagrams associated with the various \(S'_n\to S_n\). By the
diagonalizable case already proved, \(u'\) and \(u''\) are bijective.
Thus \(u\) is bijective as well, completing the proof of
\SGARefLink{sga3:ix:theorem:7:1}{7.1}. \qed

\medskip
\noindent\SGARefTarget{sga3:ix:corollary:7:3}\textbf{Corollary 7.3.}
\textit{Under the hypotheses of
\SGARefLink{sga3:ix:theorem:7:1}{7.1}, suppose in addition that \(G\)
is smooth over \(S\), and let
\[
  u_0:H_0\longrightarrow G_0
\]
be a homomorphism of \(S_0\)-groups. Then there is a homomorphism of
\(S\)-groups
\[
  u:H\longrightarrow G
\]
lifting \(u_0\). Any two such liftings \(u,u'\) are conjugate by an
element \(g\in G(S)\) whose reduction is the unit element of
\(G_0(S_0)\).}

\emph{Proof.}
This follows by combining
\SGARefLink{sga3:ix:theorem:3:6}{3.6} and
\SGARefLink{sga3:ix:theorem:7:1}{7.1}. To construct \(u\), construct
the \(u_n\) successively, as is possible by \SGARefLink{sga3:ix:theorem:3:6}{3.6}; \SGARefLink{sga3:ix:theorem:7:1}{Theorem 7.1} then shows
that the system \((u_n)\) comes from an \(S\)-homomorphism \(u\). Given
two liftings \(u,u'\), construct successively elements \(g_n\) such
that
\[
  g_0=1,
  \qquad
  g_n\text{ is the reduction of }g_{n+1},
  \qquad
  u'_n=\operatorname{int}(g_n)u_n.
\]
Again, this is possible by \SGARefLink{sga3:ix:theorem:3:6}{3.6}. Since \(A\) is separated and complete,
the \(g_n\) come from an element \(g\in G(S)\). Finally, the injectivity
in \SGARefLink{sga3:ix:theorem:7:1}{7.1} gives
\[
  u'=\operatorname{int}(g)u.
\]
\qed

\medskip
\noindent\SGARefTarget{sga3:ix:remark:7:4}\textbf{Remark 7.4.}
\textit{Compare \SGARefLink{sga3:ix:theorem:7:1}{7.1} with EGA III,
5.4.1, which implies that the statement of \SGARefLink{sga3:ix:theorem:7:1}{7.1} remains true if, instead
of assuming \(H\) of multiplicative type and \(G\) affine, one assumes
that \(H\) is proper over \(S\), and that \(G\) is separated and locally
of finite type over \(S\). A statement such as \SGARefLink{sga3:ix:theorem:7:1}{7.1} without a properness
hypothesis is rather exceptional, and should be understood here as one
aspect of the great ``rigidity'' of the structure of a group of
multiplicative type.}

\textit{The analogous statement with
\[
  G=H=\mathbf G_a
\]
is false in general. Indeed, take \(A\) of characteristic \(p>0\), and
define the \(u_n\), by reduction modulo \(I^{n+1}\), from an additive
formal power series
\SGARefTarget{sga3:ix:remark:7:4:display:01}
\[
  u(T)=\sum_{n\geq0}a_nT^{p^n},
\]
where the \(a_n\in A\) tend to zero for the \(I\)-adic topology, but not
for the discrete topology.}

\textit{The statement of \SGARefLink{sga3:ix:theorem:7:1}{7.1} also becomes false if the hypothesis that
\(G\) be affine is omitted, even when \(H=\mathbf G_m\). An example,
with \(A\) a complete discrete valuation ring, is obtained from an
elliptic curve over the fraction field \(K\) of \(A\) whose reduction,
in the N\'eron--Kodaira theory, say, is the group \(\mathbf G_m\) over the
residue field \(k\). One then has a smooth commutative group scheme \(G\)
over \(S\), with special fiber \(\mathbf G_{m,k}\). By \SGARefLink{sga3:ix:theorem:3:6}{3.6} this permits
the construction of a projective system of isomorphisms
\[
  u_n:H_n\xrightarrow{\ \sim\ }G_n,
  \qquad H=\mathbf G_{m,A},
\]
while the generic fiber of \(G\) is an abelian variety. Consequently,
there is no homomorphism of \(S\)-groups \(H\to G\) other than the zero
homomorphism.}

